Napjaink felgyorsult világában kevesen használják a tömegközlekedést, mert az nem elég kényelmes, és nem elég gyors. A buszok általában nem tartják a menetrendet, és a jegyek megvásárlása is nehézkes. Rosszabb esetben a menetrend sincs kifüggesztve a megállókban és az interneten sem elérhető. Az is gyakran megesik hogy a buszok nagyon zsúfoltak és erről az utazni vágyók csak akkor vesznek tudomást, amikor már a megállóban vannak és nincs más utazási lehetőség. A Bus4U webalkalmazás ezen problémákra kínál megoldást.

Az alkalmazás lehetőséget biztosít online jegyvásárlásra, illetve megtekinthetőek a járatok indulási időpontjai városokra és megállókra lebontva. Ezen kívül megtekinthetőek a buszmegállók a térképen, városok és útvonalak szerint szűrve, azért, hogy a felhasználó ki tudja választani a legoptimálisabb helyet a felszálláshoz. A live map funkcióban valós időben követhetőek az autóbuszok pillanatnyi helyzetei és telítettségei, hogy a felhasználó előre informálódni tudjon a menetrend változásáról és az esetleges zsúfoltságról. Ezen funkciók biztosításához, illetve a jegykezeléshez, szükség van egy alkalmazásra, ami a buszokon található POS terminálon fut. Ez biztosítja a jegyek érvényesítését, illetve az autóbuszok élő információinak megosztását a központi szerverrel.

Ezen felül a webalkalmazás biztosít egy adminisztrációs felületet is a busztársaságok számára. Ebben lehetőség nyílik elsősorban a menetrendek kezelésére, megállók és útvonalak feltöltésére. Itt lehet kezelni az alkalmazottak információit, és az autóbuszok különböző adatait.

Ez a dolgozat a webalkalmazás felhasználói felületére fókuszál, a backend működése egy másik dolgozat keretein belül van ismertetve.